\section{The contract-DAG compiler}

This section is the main compositional result.  It packages a long proof into
local nodes and reduces global correctness to a single bounded trace
invariant.

\subsection{Local certified transducers}

\begin{definition}[Contract node]
A contract node \(v\) has:
\begin{enumerate}[label=(\alph*)]
\item a fixed list of parent nodes \(\mathrm{par}(v)\);
\item a request code \(r_v\);
\item for each parent \(u\), a monotone demand transformer
      \(D_{v,u}(r_v)\);
\item a primitive-recursive constructor or finite candidate-list constructor
      \(G_v\);
\item a primitive-recursive Boolean verifier \(V_v\);
\item an internally provable local soundness sequent saying that verified
      parent observations at the demanded requests imply verification of the
      node output.
\end{enumerate}
All arities and the graph itself are fixed.  Ordinary numerical parameters
remain inputs.
\end{definition}

A local theorem may use several premises, but those premises are finite
equalities or Boolean checks.  The node is therefore a small proof-producing
transducer.  Examples include rational addition, a Taylor remainder step, one
bisection update, polynomial deflation under a nonzero leading coefficient,
and verification of one chart in a finite cover.

\subsection{Reverse demand propagation}

Let the nodes be numbered topologically \(0,\ldots,N-1\), so every parent of
\(v\) has a smaller number.  Suppose the sinks receive external requests
\(q_v\), with \(q_v=0\) for unrequested sinks.  Define required requests
backward:
\[
  r_u =
  \max\left(
       q_u,\;
       \max_{\substack{v>u\\u\in\mathrm{par}(v)}}
       D_{v,u}(r_v)
      \right).
\]
Because \(N\) and the edge relation are fixed, this is a primitive-recursive
finite reverse loop.  Monotonicity of the demand transformers guarantees that
over-approximating a downstream request never invalidates an upstream
observation.

\begin{lemma}[Demand adequacy]
For every edge \(u\to v\),
\[
        r_u\ge D_{v,u}(r_v).
\]
\end{lemma}

\begin{proof}
The right-hand side is one of the entries included in the finite maximum
defining \(r_u\).  Internally, the maximum is a primitive-recursive fold and
the claim follows from its standard quantifier-free invariant:
after processing the first \(j\) outgoing edges, the accumulator dominates
their \(j\) demands.
\end{proof}

\subsection{The forward trace}

The evaluator stores a table \(T_t\) after processing the first \(t\) nodes.
At node \(t\), it reads the already stored parent observations, applies
\(G_t\) at request \(r_t\), and appends the output.  Relational nodes append a
finite candidate list and, when necessary, a finite compatibility witness.
The trace has exactly \(N\) stages; it is not an unbounded search.

Define \(\Inv(t,T_t)\) to say:
\begin{enumerate}[label=(\roman*)]
\item \(T_t\) has the expected length and every stored code is within its
      node's primitive-recursive bound;
\item for every \(v<t\), \(V_v\) accepts the stored output at request \(r_v\);
\item every parent-child compatibility check whose endpoints are below \(t\)
      succeeds;
\item the table entries used by later nodes are the entries produced by the
      indicated earlier nodes.
\end{enumerate}
Each universal phrase is a Boolean fold over a range bounded by \(N\) or by a
stored finite-list length.  Hence \(\Inv\) is represented by a
primitive-recursive Boolean term and the formula \(\Inv=1\) is
quantifier-free.

\begin{theorem}[Finite-trace reflection]
Assume an internally proved base sequent
\[
      \mathsf{InputGuard}(d)=1 \;\vdash\; \Inv(0,T_0)=1
\]
and, uniformly in \(t<N\), an internally proved step sequent
\[
      \Inv(t,T_t)=1
      \;\vdash\;
      \Inv(t+1,T_{t+1})=1.
\]
Then iPRA proves \(\Inv(N,T_N)=1\).  If the final verifier is a component of
\(\Inv(N,T_N)\), iPRA proves the requested output contract.
\end{theorem}

\begin{proof}
Totalize the step after stage \(N\) by leaving the state unchanged, and
strengthen the invariant to
\[
   J(t,T)=
   \bigl(t\le N\wedge \Inv(t,T)=1\bigr)
   \;\vee\;
   \bigl(N<t\wedge T=T_N\wedge\Inv(N,T_N)=1\bigr).
\]
All comparisons and occurrences of \(\Inv\) are primitive-recursive Boolean
equalities, so the matrix is quantifier-free.  The base proof gives \(J(0)\).
For the successor step, decide the Boolean \(t<N\).  In the first branch use
the supplied local step theorem; in the second use the frozen-state defining
equation.  Applying quantifier-free induction yields \(\forall t\,J(t)\).
Instantiate at \(N\) and project the desired conjunct.  No induction formula
quantifies over future precisions or over arbitrary functions.
\end{proof}

\begin{theorem}[Contract-DAG composition]
For a fixed finite acyclic graph of contract nodes, suppose:
\begin{enumerate}[label=(\arabic*)]
\item every numerical operation, demand map, constructor, bound, and verifier
      is represented by a fixed \code{PR} code;
\item every local node soundness theorem is supplied as a
      \(\Prf\,\PRAAx\) derivation;
\item the input guards and source presentations are internally certified.
\end{enumerate}
Then there is a fixed primitive-recursive global evaluator, a fixed
quantifier-free global verifier, and a \(\Prf\,\PRAAx\) derivation of the
global verifier for every requested sink precision.
\end{theorem}

\begin{proof}
First run the fixed reverse loop computing the request table.  Demand adequacy
is proved by a quantifier-free fold invariant.  Next run the fixed forward
trace.  Its base case consists of the source-presentation contracts.  At step
\(v\), the trace invariant supplies the verified parent observations;
demand adequacy supplies enough precision; substitution into the local
soundness derivation yields \(V_v=1\).  Table length, bounds, and compatibility
are preserved by the defining equations of append and lookup.  The
finite-trace reflection theorem proves the final invariant.

All proof composition is intuitionistic: conjunction introduction collects
verified premises, implication elimination applies a local contract, and
bounded case splits are performed on explicit Boolean equalities.  The
external request remains a free number throughout the induction.  Universal
introduction over that request is performed only after the finite trace proof.
\end{proof}

\subsection{Why a DAG rather than a monolithic proof}

The graph exposes precisely where a packet fails.  A node may be missing:
a constructor, a numerical modulus, a separation guard, a local arithmetic
lemma, or an object-language proof.  The rest of the graph remains reusable.
It also permits shared subcomputations and shared error budgets.  Most
importantly, it gives a single formalization target:
\[
       \code{PRA.contractDAGReflection}.
\]
Once that theorem and a library of node contracts are available, future
packets should assemble proof objects rather than reproduce a bespoke
induction argument.
